\section*{ID7099: Summary / Next Steps: A(m)}
\paragraph{Metadata.}
PiID: 4363465537949 | RTEID: RTE:P27566.9935:T1.4396 | UUID: a754f604-00cf-50e4-984c-6e2b9816b914

\paragraph{Canonical Statement.}
\# 10) Summary / next steps - **Model:** \textbackslash{}(a(m)=g\_0(1+\textbackslash{}alpha m/M)\textbackslash{}), measurable difference \textbackslash{}(\textbackslash{}Delta a=g\_0\textbackslash{}alpha\textbackslash{}Delta m/M\textbackslash{}). - **Key obstacle:** Earth-based experiment with Earth as the source demands detection sensitivities \textasciitilde{}\textbackslash{}(10\textasciicircum{}\{-24\}\textbackslash{}) m/s² for α≈1 with Δm≈1 kg — effectively impossible now. - **Best bet:** perform a **local source mass experiment** (smaller \textbackslash{}(M\_s\textbackslash{}) but large local field), or maximize Δm, and use differential techniques (lock-in modulation, torsion balance, atom interferometry) to amplify SNR. Convert measured Δa to α via \textbackslash{}(\textbackslash{}alpha=(\textbackslash{}Delta a/g\_0)(M/\textbackslash{}Delta m)\textbackslash{}). - **If you want**, I can: - (A) produce a concrete experimental design using a local source mass with numbers (choose source mass, separation, geometry), and compute expected \textbackslash{}(g\_s\textbackslash{}) and Δa for a range of α, or - (B) give a template data-analysis script (Python) that converts measured Δa ± σ into α ± σ and plots confidence intervals, - (C) show how to interpret atom-interferometer results if your α couples per-nucleon rather than total mass.

\paragraph{Canonical Equation.}
\[
a(m)=g_0(1+\alpha m/M)
\]

\paragraph{Dictionary.}
Relation: a(m)=g\_0(1+α m/M)

\paragraph{Encyclopedia.}
Summary / Next Steps: A(m) is an algebraic definition, written a(m)=g\_0(1+α m/M). It is an algebraic definition expressing the quantity directly in terms of the variables on its right-hand side. It is built from α, g\_0, m, M, defining a(m). Within the Trawin Zero Point registry it serves as one element of the framework's mathematical narrative.
